% ============================================================
% Section 79: 抛物线型周期势的能带与有效质量
% ============================================================
\section{固体理论：抛物线型周期势的能带与有效质量}
\label{sec:parabolic-potential-band}

\begin{modelbox}[题目：抛物线型周期势的近自由电子能带与有效质量]
电子周期场的势能函数为
\[
V(x)=\frac{1}{2}m\omega^2\bigl[b^2-(x-na)^2\bigr],\quad na-b\le x\le na+b,
\]
其中 $a=4b$，$\omega$ 为常数。
\begin{enumerate}
    \item 画出势能曲线，并求其平均值；
    \item 用近自由电子近似，求第一和第二个禁带宽度；
    \item 若每个原胞含 4 个电子，该晶体是金属还是非金属？
    \item 求价带顶及导带底处的电子有效质量；
    \item 求价带顶及导带底处的电子平均速度。
\end{enumerate}
\end{modelbox}

\begin{tipbox}[考点快速分析]
\begin{itemize}
    \item \textbf{抛物线势}：倒置抛物线，在 $x=na$ 处取最大值 $V_{\max}=\frac12m\omega^2b^2$
    \item \textbf{傅里叶积分}：$V_n=\frac{1}{a}\int_{-b}^{b}V(x)e^{-i2\pi nx/a}dx$，利用分部积分
    \item \textbf{有效质量}：$m^*=\hbar^2\bigl(\frac{\dd^2E}{\dd k^2}\bigr)^{-1}$；带底 $m^*>0$，带顶 $m^*<0$
    \item \textbf{带边速度}：$v=\frac{1}{\hbar}\frac{\dd E}{\dd k}$，极值点处为零
\end{itemize}
\end{tipbox}

\begin{derivationbox}[标准解题过程]
\textbf{Step 1: 势能曲线与平均值}

势能是以 $a=4b$ 为周期的周期函数。在一个周期内（$n=0$，$-b\le x\le b$），
\begin{equation}
V(x)=\frac12m\omega^2(b^2-x^2).
\end{equation}
在 $x=0$ 处取最大值 $V_{\max}=\frac12m\omega^2b^2$，在 $x=\pm b$ 处 $V=0$。相邻势垒之间有宽度为 $2b$ 的零势区（因 $a=4b$）。

平均值（在一个周期 $[-2b,2b]$ 内积分）：
\begin{align}
V_0&=\frac{1}{4b}\int_{-2b}^{2b}V(x)\,dx=\frac{1}{4b}\int_{-b}^{b}\frac12m\omega^2(b^2-x^2)\,dx \notag\\
&=\frac{m\omega^2}{8b}\Bigl[b^2x-\frac{x^3}{3}\Bigr]_{-b}^{b}
=\frac{m\omega^2}{8b}\cdot\frac{4b^3}{3}=\boxed{\frac{m\omega^2b^2}{6}.}
\end{align}

\textbf{Step 2: 第一和第二禁带宽度}

傅里叶系数
\begin{equation}
V_n=\frac{1}{a}\int_{-a/2}^{a/2}V(x)e^{-i\frac{2\pi n}{a}x}\,dx
=\frac{1}{4b}\int_{-b}^{b}\frac12m\omega^2(b^2-x^2)e^{-i\frac{n\pi}{2b}x}\,dx.
\end{equation}

利用偶函数性质，仅余弦项贡献。经分部积分计算（过程略），结果为
\begin{equation}
V_n=-\frac{2m\omega^2b^2}{n^2\pi^2}\cos\frac{n\pi}{2}+\frac{4m\omega^2b^2}{n^3\pi^3}\sin\frac{n\pi}{2}.
\end{equation}

\textbf{第一禁带}（$n=1$）：$\cos(\pi/2)=0$，$\sin(\pi/2)=1$
\begin{equation}
V_1=\frac{4m\omega^2b^2}{\pi^3}
\quad\Longrightarrow\quad
\boxed{E_{g,1}=2|V_1|=\frac{8m\omega^2b^2}{\pi^3}.}
\end{equation}

\textbf{第二禁带}（$n=2$）：$\cos\pi=-1$，$\sin\pi=0$
\begin{equation}
V_2=\frac{2m\omega^2b^2}{4\pi^2}=\frac{m\omega^2b^2}{2\pi^2}
\quad\Longrightarrow\quad
\boxed{E_{g,2}=2|V_2|=\frac{m\omega^2b^2}{\pi^2}.}
\end{equation}

\textbf{Step 3: 每原胞 4 个电子的导电性}

一维晶格每条能带可容纳 $2N$ 个电子（$N$ 为原胞数）。每原胞 4 个电子，总电子数 $4N$：
\begin{itemize}
    \item 第一能带填满（$2N$ 个电子）
    \item 第二能带填满（$2N$ 个电子）
\end{itemize}
由于第二禁带 $E_{g,2}\neq0$，第二能带与第三能带之间存在非零能隙。在绝对零度下，前两带全满，第三带全空，费米能级落在禁带中。因此材料为\textbf{非金属（绝缘体）}。

\textbf{Step 4: 价带顶与导带底的有效质量}

在近自由电子模型中，布里渊区边界 $k=G/2=n\pi/a$ 附近，能带可近似为
\begin{equation}
E_{\pm}(k')=E_0\pm|V_n|+\frac{\hbar^2k'^2}{2m}\Bigl(1\pm\frac{2E_0}{|V_n|}\Bigr)^{-1},
\end{equation}
其中 $k'$ 是从边界点测量的波矢，$E_0=\hbar^2G^2/(8m)=\hbar^2n^2\pi^2/(2ma^2)$ 为边界处自由电子能量。

有效质量定义为
\begin{equation}
\frac{1}{m^*}=\frac{1}{\hbar^2}\frac{\dd^2E}{\dd k'^2}=\frac{1}{m}\Bigl(1\pm\frac{2E_0}{|V_n|}\Bigr)^{-1}.
\end{equation}

对于第二布里渊区边界（$n=2$，$G=4\pi/a=\pi/b$）：
\begin{equation}
E_0=\frac{\hbar^2G^2}{8m}=\frac{\hbar^2\pi^2}{8mb^2},\qquad |V_2|=\frac{m\omega^2b^2}{2\pi^2}.
\end{equation}

比值
\begin{equation}
\frac{2E_0}{|V_2|}=\frac{\hbar^2\pi^4}{2m^2\omega^2b^4}.
\end{equation}

\textbf{导带底}（上支，$+$号）：
\begin{equation}
\boxed{m_c^*=\frac{m}{1+\dfrac{\hbar^2\pi^4}{2m^2\omega^2b^4}}\approx\frac{2m^3\omega^2b^4}{\hbar^2\pi^4}\;(2E_0\gg|V_2|\text{时}).}
\end{equation}

\textbf{价带顶}（下支，$-$号）：
\begin{equation}
\boxed{m_v^*=\frac{m}{1-\dfrac{\hbar^2\pi^4}{2m^2\omega^2b^4}}\approx-\frac{2m^3\omega^2b^4}{\hbar^2\pi^4}\;(2E_0\gg|V_2|\text{时}).}
\end{equation}
导带底有效质量为正（电子型），价带顶有效质量为负（空穴型）。在近自由电子极限（$|V_2|\ll E_0$），$|m^*|\ll m$。

\textbf{Step 5: 价带顶与导带底的电子平均速度}

电子平均速度（群速度）
\begin{equation}
v=\frac{1}{\hbar}\frac{\dd E}{\dd k'}.
\end{equation}
在能带极值点（带顶或带底），$\dd E/\dd k'=0$，因此
\begin{equation}
\boxed{v_c(k'=0)=0,\qquad v_v(k'=0)=0.}
\end{equation}

物理上，带边处电子形成驻波态（布洛赫波在布里渊区边界发生布拉格反射），净传播速度为零。
\end{derivationbox}

\begin{formulabox}[答案汇总]
\begin{center}
\begin{tabular}{ll}
\toprule
\textbf{问题} & \textbf{答案} \\
\midrule
(1) 势能平均值 & $V_0=m\omega^2b^2/6$ \\
(2) 第一禁带宽度 & $E_{g,1}=8m\omega^2b^2/\pi^3$ \\
(2) 第二禁带宽度 & $E_{g,2}=m\omega^2b^2/\pi^2$ \\
(3) 导电性 & 4电子/原胞填满两带，存在非零禁带；\textbf{绝缘体} \\
(4) 导带底有效质量 & $m_c^*=m\bigl/\bigl(1+\hbar^2\pi^4/(2m^2\omega^2b^4)\bigr)$ \\
(4) 价带顶有效质量 & $m_v^*=m\bigl/\bigl(1-\hbar^2\pi^4/(2m^2\omega^2b^4)\bigr)$ \\
(5) 带边平均速度 & $v_c=v_v=0$ \\
\bottomrule
\end{tabular}
\end{center}
\end{formulabox}

\begin{insightbox}[物理图像]
\begin{itemize}
    \item 抛物线势的傅里叶分量同时含正弦和余弦贡献，奇偶阶均非零，故各阶禁带均可能存在（与方波势不同）。
    \item 有效质量反映了能带曲率。在弱周期势下，能带在边界处非常``平坦''（曲率大），导致有效质量很小。这相当于电子在晶格中运动时，周期势的布拉格反射强烈地改变了电子的惯性响应。
    \item 带边速度为零是布洛赫波驻波特性的直接体现。电子在带边处不是传播波，而是被晶格周期势``锁定''的驻波。
    \item 价带顶有效质量为负是半导体物理中``空穴''概念的来源：用正电荷、正有效质量的空穴描述价带顶附近电子的集体行为，比直接处理大量负有效质量电子更方便。
\end{itemize}
\end{insightbox}
